\section{Visão geral do projeto}

O projeto implementa um simulador de resposta a surto em bobina distribuída. A ideia central é substituir o enrolamento contínuo por uma cadeia de seções concentradas, cada uma contendo resistência, indutância e capacitância para terra. Essa discretização permite estudar fenômenos de alta frequência que um modelo puramente concentrado perderia, como atraso de propagação, reflexões na extremidade aberta e distribuição não uniforme de tensão.

A pasta \codefile{config} contém o caso base em \codefile{default\_case.json}. A pasta \codefile{src} contém o modelo elétrico, a fonte, o solver, o processamento de resultados e a visualização. A pasta \codefile{output} contém os CSVs, PNGs e GIFs já gerados. O arquivo \codefile{surto\_bobina.atp} descreve um caso ATP com a mesma lógica física do modelo Pi: \Nsec{} ramos série R--L, capacitores shunt para terra e uma fonte dupla exponencial aplicada no nó de entrada.

Os resultados numéricos são exportados como tensões nodais, correntes de seção e um resumo com picos de tensão e gradientes. Os resultados visuais incluem curvas de entrada/saída, curvas em nós selecionados, mapas de calor posição-tempo, barras de pico por nó, barras de gradiente entre nós adjacentes e animações da onda viajante.

\begin{table}[H]
\centering
\caption{Parâmetros principais do caso padrão em \codefile{config/default\_case.json}.}
\label{tab:parametros-caso}
\begin{tabular}{lll}
\toprule
Parâmetro & Valor & Interpretação elétrica \\
\midrule
\codefile{n\_sections} & 20 & número de seções discretas \\
\codefile{L\_total} & \SI{0.01}{H} & indutância total da bobina \\
\codefile{R\_total} & \SI{5}{\ohm} & perdas série totais \\
\codefile{C\_total} & \SI{1e-9}{F} & capacitância total para terra \\
\codefile{model\_type} & \codefile{pi} & topologia Pi principal \\
\codefile{source\_type} & \codefile{double\_exp} & impulso dupla exponencial \\
\codefile{V\_amplitude} & \SI{1000}{V} & pico nominal da fonte \\
\codefile{t\_front} & \SI{1.2}{\micro\second} & tempo de frente \\
\codefile{t\_tail} & \SI{50}{\micro\second} & tempo de cauda \\
\codefile{dt} & \SI{10}{ns} & passo de amostragem dos resultados \\
\codefile{termination} & \codefile{open} & extremidade aberta \\
\bottomrule
\end{tabular}
\end{table}
